\documentclass[10pt]{article}
\usepackage[margin=20mm]{geometry}
\usepackage{amsmath,amssymb,mathtools,fontspec,unicode-math,xurl,hyperref}
\IfFontExistsTF{DejaVu Serif}{\setmainfont{DejaVu Serif}}{\setmainfont{Latin Modern Roman}}
\setmathfont{latinmodern-math.otf}
\setlength{\parindent}{0pt}\setlength{\parskip}{6pt}\sloppy
\title{Technical Review of Plane Quartic
Reduction via Cayley Octads}
\author{AI-assisted technical review}
\date{}
\begin{document}
\maketitle
Abstract: A focused review finds a Minor scope defect in Proposition 6.3: the stated classification of all merotropic subspaces omits a four-dimensional coisotropic example. The classification is verified after restricting it to the smaller-dimensional representatives used in the subsequent compatibility definition. Exact determinant and finite-field checks support the inspected examples and Table 8. The geometric classification remains explicitly conjectural in the paper, and the large computational proofs and database experiments were not independently reproduced. Scientific adjudication is pending.\par
\section*{Identity, scope and logical dependencies}
The reviewed paper is R. van Bommel, J. Docking, V. Dokchitser, R. Lercier and E. Lorenzo García, Reduction of Plane Quartics and Cayley Octads, Foundations of Computational Mathematics 26 (2026), 1425–1496, DOI https://doi.org/10.1007/s10208-025-09704-y. The exact source is the 72-page publisher version of record, published online 2 June 2025. Physical PDF page \(p\) has printed page \(1424+p\). Locations below use physical pages unless specified.\par
For a plane quartic over a non-archimedean local field with residue characteristic different from two, the paper associates valuation data to a Cayley octad of eight points in projective three-space. After normalizing each point to have minimum coordinate valuation zero, the 70 Plücker coordinates are the four-column determinants. A further twisted-cubic index records additional degeneration. Compatible building blocks give an octad picture.\par
The distinction between claims is essential. Existence of compatible decompositions is Conjecture 5.4 (p. 42); their geometric interpretation as stable reduction is Conjecture 5.12 (p. 44) and Conjecture 6.27 (p. 52). General symplectic equivariance is Conjecture 6.13 (p. 47). Theorems 5.5 and 5.8 concern uniqueness and coordinate invariance in the compatible-decomposition framework. Theorem 6.29 (p. 53) is a finite combinatorial graph correspondence, not an unconditional proof of the geometric conjecture. The absence of a general geometric proof is not reported as a defect.\par
The review covers the main scope statements, determinant examples, valuation-normalization argument, selected finite-field combinatorics, static computational-proof scripts, and the stated evidence and family-construction limitations. It does not certify all 42 stable types or all appendix correspondences.\par
\section*{Problem 1: the subspace classification needs the smaller-side restriction}
Location: Definition 6.2 and Proposition 6.3, p. 45 (printed p. 1469), followed by Table 8 and Definition 6.4. Assessment: Minor.\par
Definition 1.19 realizes \(E_3\) as the even subsets of an eight-letter alphabet, modulo complementation. Addition is symmetric difference and \[\langle I,J\rangle=|I\cap J|\pmod2.\] This is a six-dimensional symplectic space over \(\mathbb F_2\). Definition 6.2 calls a subspace merotropic unless it is isotropic of dimension greater than one. Proposition 6.3 claims that, up to \(S_8\) and orthogonal complementation, all such subspaces have one of the eight listed representatives.\par
Set \[Y=\langle AB,CD\rangle,\qquad X=Y^\perp.\] The two independent generators of \(Y\) are orthogonal, so \(Y\) is isotropic of dimension two. Nondegeneracy of the ambient pairing gives \[\dim X=4,\qquad X^\perp=Y,\qquad
 X\cap X^\perp=Y,\qquad
 \operatorname{rank}(\langle\, ,\,\rangle|_X)=4-2=2.\] Thus \(X\) is not isotropic and is merotropic under the stated definition.\par
Neither \(X\) nor its complement can occur in the printed list. Dimensions rule out its zero-, one-, and three-dimensional representatives. The two listed two-dimensional representatives are \(\langle AB,AC\rangle\) and \(\langle AB,ACDE\rangle\). Their generators have pairing one, so both are nondegenerate. Their complements have dimension four and rank four. The unordered dimension/rank signature of our pair is \[\{(2,0),(4,2)\},\] whereas those listed two/four-dimensional pairs have signature \(\{(2,2),(4,4)\}\). Permuting letters preserves the pairing and hence cannot eliminate this distinction. All assumptions on \(X\) have been verified; no geometric realization by a quartic is required for this finite-vector-space proposition.\par
There is a natural narrower reading. Definition 6.4 orders a pair \((X,X^\perp)\) so that \(\dim X\le\dim X^\perp\), and then requires that this smaller member be merotropic. Our example is excluded by that requirement because its smaller member is \(Y\). A precise correction to Proposition 6.3 is therefore: classify orthogonal pairs whose smaller-dimensional member is merotropic, with either representative allowed when dimensions are equal. This is a genuine restriction on the classified objects, rather than a change to the pairing.\par
The corrected claim was independently checked by exhaustive finite enumeration from the subset definition. Dimensions zero through three contain \(1,63,651,1395\) subspaces. After removing isotropic subspaces of dimensions two and three, the counts are \(1,63,336,1260\). Quotienting the three-dimensional case by orthogonal complementation and taking \(S_8\)-orbits gives: \[\begin{array}{c|l}
\text{smaller dimension}&\text{orbit sizes}\\\hline
0&1\\
1&28,\ 35\\
2&56,\ 280\\
3&35,\ 280,\ 315
\end{array}\] There are exactly eight orbits. Direct span and pairing calculations also verify every orthogonal-complement and radical entry of Table 8. The defect is local to the unrestricted wording; the actual Definition 6.4 domain is unchanged by this repair. No failure of the compatible-subspace machinery or stable-reduction conjecture follows. That bounded impact supports Minor, not a core-contribution Major.\par
\section*{Exact checks of valuation data and normalization}
\subsection*{Examples 1.1, 1.9 and 3.8}
For Example 1.1 (pp. 4–5), every four-column determinant of the displayed rational octad is nonzero. Reduction modulo 17 confirms that its first four columns lie on \(s=0\) and the last four on \(6u+9v+2r=0\). This verifies the stated incidence data, not the subsequent contraction or stable-model identification.\par
For Example 1.9 (p. 9), all 70 determinants were computed exactly. Their 7-adic valuation histogram is \[\#\{v=0\}=54,\qquad\#\{v=1\}=15,\qquad\#\{v=2\}=1.\] For every four-element subset \(S\), \[v_S=\max(0,|S\cap\{A,F\}|-1)
      +\mathbf1_{S=\{A,E,F,H\}}
      +\mathbf1_{S=\{B,C,D,G\}}.\] Thus the Plücker part of the decomposition in Example 1.12 (p. 11) is exact: the \(AEFH\) coordinate is two, the \(BCDG\) coordinate is one, and the remaining positive coordinates reflect the \(AF\) collision. All \(8!=40320\) permuted instances of the Cayley relation (2), p. 27, also vanish exactly for this matrix. These include redundant instances and are not claimed to be 40320 independent relations.\par
Example 3.8 (p. 30) yields the stated histogram \(16,48,6\) at valuations \(0,1,2\) for \(p=10007\). Every coordinate equals the sum of the four point-coincidence vectors for \(AE,BF,CG,DH\). The 840 twisted-cubic relations, the generic comparison used in the twisted-cubic index, and the geometry of the corresponding stable fibres were not reproduced.\par
\subsection*{The determinant update and the normalization coset}
In Proposition 3.13 (pp. 31–32), let \(S\) index precisely the points reducing to \((1:0:0:0)\) and let \(T=S^c\). Multiplication by \(\operatorname{diag}(\pi,1,1,1)\) adds one to every four-column determinant valuation. Renormalizing a point in \(S\) subtracts one from each determinant containing it. For a quadruple \(Q\), the net change is therefore \[1-|Q\cap S|
 =-\max(0,|Q\cap S|-1)+\max(0,|Q\cap T|-3),\] exactly the claimed point/plane update, including the cases \(|Q\cap S|=0\) and \(4\). Twisted-cubic invariance uses the balanced-binomial premise of Proposition 3.12; its external 840-polynomial input was not independently reconstructed.\par
The normalization uniqueness behind Proposition 3.15 can be seen without relying on the repeated coordinate labels in its displayed proof. A determinant valuation changes by \[w_Q=d+\sum_{i\in Q}c_i=\sum_{i\in Q}a_i,
 \qquad a_i=c_i+d/4.\] The eight incidence vectors \(Q\mapsto\mathbf1_{i\in Q}\) span an eight-dimensional space \(W_8\), as verified by exact rank. Imposing zero on the five quadruples from \(A,\ldots,E\) determines \(a_A,\ldots,a_E\) uniquely: the homogeneous equations force these five numbers equal, and then four times their common value is zero. The remaining three parameters independently shift the four coordinates containing \(F\), \(G\), or \(H\) and three of \(A,B,C,D\). Setting each of these three minima to zero fixes them uniquely. Thus the normalized Plücker valuation representative depends only on the original valuation vector.\par
The point-coincidence vectors have rational rank 28, the complementary-plane vectors rank 35, and their combined rational span rank 42, consistent with Remark 3.10. Equality of integral lattices, rather than rational spans, was not verified.\par
\section*{Computational evidence and remaining proof obligations}
The acquired repository is the paper-cited commit `961e584b694a217b96e9168f16ba4a9a525c6149`. Its README and the complete scripts `UniqueDecomposition.m` and `CheckMinimaCompatibleSums.m` were inspected statically. The latter compares the sum with its full normalized valuation vector. It therefore tests all normalization requirements, rather than only a weak reading of the three-minima sentence on p. 44. For nonnegative block vectors, an unweighted sum having minimum zero means that a common coordinate is zero in every summand; all positive weighted sums then have the same property. This explains why checking supports is relevant for arbitrary positive block coefficients.\par
The uniqueness script constructs compatible subsets, reduces by letter permutations, and examines kernel signs in the quotient by \(W_8\). Its comments about 2,626,648 sets and 507 representatives are author assertions, not new runtime results. The complete generator implementation, exhaustiveness of those enumerations, sign-branch sufficiency and the reported absence of printed exceptional pairs remain unverified here. The paper’s 16-CPU-hour computation was not rerun, and no conclusion of Theorem 5.5 is certified by merely locating its code.\par
For the elementary two-line example on pp. 17–19, \(\langle AF\rangle\) and \(\langle AEFH\rangle\) are orthogonal. The displayed graph construction doubles each edge, assigns the two leaves genus zero, contracts them, and leaves a genus-one vertex with two loops. Its total genus is \(1+2=3\), as in Figure 8 and the corresponding Table 11 row. The more complicated relation-index construction and Theorem 6.29’s complete graph isomorphism were not independently verified.\par
Section 1.4 reports checks for about 180,000 of 186,502 curve/prime pairs, and explicitly describes limitations of the independent comparison for the remaining cases. Section 7.1 gives generic family constructions, with regularity conditions and special singular cases; Figure 20 distinguishes those cases. The omitted blow-up and base-change details for the three-cusp family on p. 59 were not reconstructed. Appendix A explicitly conditions its geometric use on the conjectures. Selected table pages were visually inspected, without claiming a complete row-by-row mathematical audit. No new empirical results or universal geometric proof are inferred.\par
A bounded primary-source lookup on 29 October 2026 confirmed the publisher identity and arXiv v2 dated 21 February 2024. No linked correction was located in that lookup; this is not an exhaustive absence claim. The arXiv mathematical text was not compared.\par
Item | Initial-review status\par
Issue judgment | Minor scope defect in Proposition 6.3.\par
Repair | Smaller-side classification verified exactly; Table 8 verified.\par
Other checks | Determinant examples and normalization argument supported within the stated scope.\par
Author adoption | Unknown; not requested.\par
Remaining work | Large uniqueness and graph enumerations, geometric families, twisted-cubic index and empirical reproduction.\par
Scientific decision | Pending; human reviews empty; this is not independent acceptance.\par
\section*{Scope and references}
PDF SHA256: `8138b720ca27fd0f1449cf5ad8c49073d60f50dc5cba670742c00ccfa03ec946`. Manifest SHA256: `8b3fd6d3951132e8d2910f7a9a9991c1048eaad37cdbb561765eaa62db2b4da2`. All 72 pages were mechanically parsed. The coverage ledger distinguishes scientific reading, partial inspection and unreviewed pages. Independent `exact\_checks.py` and `scope\_check.py` both exited zero; their complete code, commands and literal results are retained. They execute no author code. The subset enumeration uses even eight-bit masks modulo complement and adjacent letter transpositions, rather than imported author tables. Actual compiler/dependency payloads, build logs, PDF parsing and source/report image-view receipts are separate verification artifacts. All science and page inspection were performed by this assigned executor. The helper supplied clerical checklist advice only. The main dispatcher supplied the visible model/account observation; runtime self-description is separately recorded, with no backend equivalence certification.\par
9 R. van Bommel, J. Docking, V. Dokchitser, R. Lercier and E. Lorenzo García, Reduction of Plane Quartics and Cayley Octads, Found. Comput. Math. 26 (2026), 1425–1496, DOI: 10.1007/s10208-025-09704-y. Exact publisher PDF reviewed. Same authors, g3cayley, https://github.com/rbommel/g3cayley, commit `961e584b694a217b96e9168f16ba4a9a525c6149`; selected scripts read statically. Same authors and title, https://arxiv.org/abs/2309.17381v2, 21 February 2024; metadata consulted only.\par
\end{document}
